\subsection{Common Diagnostics}
\label{sec:generator-diagnostics}

Unrestricted feedback can amplify observations over repeated execution.
A pretraining prior dominated by such behaviour would allocate much
of its capacity to rapidly growing recurrences. For the intended benchmark,
explosive cases can instead form an explicitly controlled difficulty group,
alongside mechanisms producing sustained variation at usable scales.

This section defines explosion, low-amplitude, and tail-flattening diagnostics.
The experiment configuration and timing protocol are specified in
Section~\ref{sec:generator-evaluation-protocol}. Their measured trade-offs
appear in Section~\ref{sec:implemented-generator-results}, where improved
growth control is weighed against reduced variation and computational cost.

\subsubsection{Explosion Detection}

The evaluator constructs the history according to each family, generates one
path up to the largest horizon, and reuses its prefixes. No generated prefix
is discarded before these observations: the original, bounded-domain, and
Shao--Wu families draw their initial histories directly, while Heaps draws
from its joint stationary law. It stops at the first
non-finite value or magnitude greater than 100, retaining that failing value.
The event remains counted at every later horizon, including a failure at the
last requested step. Every attempted mechanism stays in the denominator.
For multivariate use, any failing channel flags the system. These rules are
implemented in \path{src/tsfmxai/evaluation/stability.py} and
\path{src/tsfmxai/evaluation/generator_tradeoff.py}.

An amplitude threshold is an operational safeguard: Gaussian innovations have
unbounded support and can cross any finite threshold even under a stationary
law. The evaluator records finite threshold crossings separately from
non-finite outputs.

\subsubsection{Near-Zero Amplitude}
A completed, non-explosive prefix receives this label when every channel's
last 24 values lie in $[-0.15,0.15]$. The tolerance is three times the
innovation standard deviation used in this experiment, $\sigma_\varepsilon=0.05$.
Thus the label measures occupation of a small band during a finite window.
It can also be produced by a stationary process fluctuating near zero from
the beginning. Convergence of a trajectory to zero is a stronger,
infinite-time property requiring separate evidence.

\subsubsection{Nearly Constant Tails}
To capture stabilization around levels other than zero, let
$\bar x^{\mathrm{tail}}_{k,h}$ be the mean of channel $k$ over the last
$w_{\mathrm{tail}}$ generated observations at horizon $h$. Define
\begin{equation}
 V_h^{\mathrm{tail}}=
 \max_k\max_{h-w_{\mathrm{tail}}<v\leq h}
 \left|x_{k,\tau+v}-\bar x^{\mathrm{tail}}_{k,h}\right|.
 \label{eq:generator-flat-tail}
\end{equation}
A completed, non-explosive prefix is nearly constant when
$V_h^{\mathrm{tail}}\leq\delta_{\mathrm{flat}}$. The implementation uses
$w_{\mathrm{tail}}=24$ and the operational tolerance
$\delta_{\mathrm{flat}}=3\sigma_\varepsilon=0.15$, retaining each channel's
mean and variation. Persistence requires both consecutive terminal windows
to pass, and all 48 terminal observations to satisfy the same tolerance
around their pooled channelwise mean. This joint check captures level changes
between individually flat windows. Variation is measured on signed channels
before taking a maximum across channels.

The first 24 generated observations, after the family's initialization,
provide a disjoint reference window. Observed flattening requires persistent
flatness, at least one channel initially above the variation tolerance, and
a reduction of at least 50\% in every such channel. The reduction compares
the pooled 48-step terminal variation with the initial 24-step variation.
Quiet initial and terminal windows are recorded separately; intermediate
behaviour remains unrestricted. A persistent flat tail is labelled away
from zero when at least one channel's final mean lies outside $[-0.15,0.15]$.

Amplitude decline towards the near-zero band uses absolute values relative
to zero: every channel must remain in the band for the final 48 steps,
at least one channel must initially exceed the band, and every such channel
must reduce its maximum absolute value by at least 50\%. These finite-window
labels can overlap. All rates retain every attempted trajectory in the
denominator; exploded prefixes are ineligible for tail measurements.
The notebook also reports eligible-only rates and sensitivity to
$\delta_{\mathrm{flat}}\in\{\sigma_\varepsilon,2\sigma_\varepsilon,
3\sigma_\varepsilon\}$. These thresholds and the reduction factor are
configurable project choices. The implementation is in
\path{src/tsfmxai/evaluation/tail_behavior.py}.
